\documentclass[10pt,a4paper]{amsart}
\usepackage[margin=19mm]{geometry}
\usepackage{amsmath,amssymb,mathtools}
\usepackage[scr=rsfs,cal=euler]{mathalfa}
\usepackage{xfrac}
\usepackage[unicode,hidelinks,bookmarks=false]{hyperref}
\newcommand{\ie}{i.e.\ }
\newcommand{\cf}{cf.\ }
\newcommand{\resp}{respectively\ }
\newcommand{\Subsection}{\subsection}
\providecommand{\nobreakdash}{\nobreak\hbox{-}}
\setlength{\emergencystretch}{2em}
\multlinegap0pt
\usepackage{tikz}
\usetikzlibrary{cd}
\usepackage{amssymb}
\usepackage{mathtools}
\usepackage{tikz}
\usepackage{cancel}
\usepackage{xcolor}
\newtheorem{theorem}{Theorem}
\newcommand{\R}{\mathbb{R}} 
\newcommand{\an}[1]{\arrowvert_{#1}} 
\newcommand{\inv}{^{-1}}
\newcommand{\mx}{\mathfrak{X}} 
\title{On the flows of linear vector fields}
\author{M. Jotz}
\date{Source: arXiv:2409.18723v3 (CC BY 4.0)}
\begin{document}
\maketitle

\noindent\emph{Source and licence.} On the flows of linear vector fields. Version arXiv:2409.18723v3 (CC BY 4.0), distributed by arXiv under the Creative Commons Attribution 4.0 International licence, http://creativecommons.org/licenses/by/4.0/, as recorded in the arXiv licence metadata for this exact version and shown on the abstract page https://arxiv.org/abs/2409.18723v3. Full licence notice: \texttt{licenses/arxiv-2409.18723-CC-BY-4.0.txt}.\par\smallskip

\noindent\textit{Editorial scope.} Counted unit: the article's theorem lemma\_linear\_flow (source0.tex lines 538--550). The article's complete original proof of it is delivered with it (source0.tex lines 553--625, one \texttt{proof} environment of 73 lines). No mathematics is written, repaired or re-derived: the statement and the proof are the article's own lines, in the article's own notation, and no retained line is substituted or re-flowed.  No further source-local context is needed: every cross-reference of every retained line resolves inside the two delivered spans.  Nothing was omitted from any retained span, so the package carries no omission record.  No retained line contains a citation, so the wrapper carries no bibliography and no bibliography entry.  The retained lines contain the article's own diagram code, which the wrapper typesets with the standard tikz packages; no image file is delivered and the package declares no asset dependency.  Editorial formatting only, in addition: an AMS \texttt{amsart} preamble that restates the article's own declarations and macros used by the retained lines, namely the environments \texttt{theorem} on separate counters, together with the standard packages those lines need.  The retained environments print in this document as Theorem 1 in order of appearance, whereas the article prints the same items with the numbering of its own full document.  The complete native source of the article as distributed in the arXiv e-print for this exact version is bundled unchanged as \texttt{sources/165931/source0.tex}.
\begin{theorem}\label{lemma_linear_flow}
Let $E\to M$ be a smooth vector bundle and let $\chi\in\mx^l(E)$ be a linear vector field on $E$ over $X\in\mx(M)$. Let $\phi\colon \Omega\to M$ be the flow of $X$, with, as usual, $\Omega\subseteq \R\times M$ an open subset containing $\{0\}\times M$. Then the flow $\Phi$ of $\chi\in\mx^l(E)$ is defined on 
\[ \widetilde\Omega:=\left\{ (t, e)\in \R\times E \mid (t,q(e))\in \Omega
\right\}
\]
and for each $t\in\R$, the map
\[ \Phi_t\colon q\inv(M_t)\to q\inv(M_{-t})
\]
is an isomorphism of vector bundles over the diffeomorphism $\phi_t\colon M_t\to M_{-t}$, where 
$M_t\subseteq M$ is the open subset
\[ M_t:=\{x\in M\mid (t,x)\in\Omega\}.
\]
\end{theorem}

\begin{proof}
First, as usual $Tq\circ \chi=X\circ q$ implies that each flow line of $\chi$ projects under $q$ to a flow line of $X$, and so that 
\[ \widetilde \Omega\subseteq \{(t, e)\in \mathbb R\times E\mid (t, q(e))\in \Omega\}.
\]

Choose a point $x\in M$ and consider the maximal integral curve $c\colon J\to M$ of $X$ through $x$ at time $0$, i.e.~with an open interval $J\subseteq \mathbb R$ containing $0$. Define $C\colon J\to E$, $C(t)=0^E(c(t))$ for all $t\in J$. Then  for all $t\in J$
\[ \dot C(t)=T_{c(t)}0^E(\dot c(t))=T_{c(t)}0^E(X(c(t)))=\chi(0^E(c(t)))=\chi(C(t))
\]
since $\chi$ is linear and so in particular 
% https://q.uiver.app/#q=WzAsNCxbMCwwLCJFIl0sWzAsMSwiTSJdLFsxLDEsIlRNIl0sWzEsMCwiVEUiXSxbMCwzLCJcXGNoaSJdLFsxLDIsIlgiLDJdLFsxLDAsIjBeRSJdLFsyLDMsIlxccXF1YWQgXFxxcXVhZCBUMF5FPTBee1RxfSIsMix7Im9mZnNldCI6MX1dXQ==
\[\begin{tikzcd}
	E & TE \\
	M & TM
	\arrow["\chi", from=1-1, to=1-2]
	\arrow["{0^E}", from=2-1, to=1-1]
	\arrow["X"', from=2-1, to=2-2]
	\arrow["{T0^E}"', shift right, from=2-2, to=1-2]
\end{tikzcd}\]
commutes. (The map $T0^E$ is the zero section of the vector bundle $Tq\colon TE\to TM$.) That is, $C\colon J\to E$ is an integral curve of $\chi$ starting at $0^E_x$ at time $0$.
This shows that if $(t,x)\in \Omega$ for $t\in\mathbb R$ and $x\in M$, then also $(t,0^E_x)\in\widetilde \Omega$, the flow domain of $\chi$.


Since $\widetilde \Omega$ is open in $\mathbb R\times E$, there exist then $I\subseteq \mathbb R$ open around $t$ and $U\subseteq E$ open around $0^E_x$ such that 
\[ (t, 0^E_x)\in I\times U\subseteq \widetilde \Omega.
\]
Since $U\subseteq E$ is open and $E$ is locally trivial, $U$ contains the elements of a basis $(e^1, \ldots, e^k)$ of $E_x$.
Then $(t,e^i)\in \widetilde \Omega$ for $i=1,\ldots, k$. Write 
\[ c^i\colon J_i\to E
\]
for the flow line of $\chi$ through $e^i$ at time $0$, hence with $0, t\in J_i$ for each $i=1,\ldots, k$. Then as already observed, the curves $q\circ c^i\colon J_i\to M$ are all integral curves of $X$ through $x$ at time $0$. They must hence coincide on the intersection
\[ I=\bigcap_{i=1}^k J_i,
\]
which is an open interval containing $0$ and $t$.
Take now an arbitrary $e\in E_x$ and write it in the basis $(e^1, \ldots, e^k)$:
\[ e=\sum_{i=1}^k \alpha_ie^i
\]
with $\alpha_1,\ldots, \alpha_k\in\mathbb R$.  
Consider the smooth curve 
\[ C\colon I\to E, \qquad s\mapsto \sum_{i=1}^k\alpha_i c^i(s).
\]
This is well-defined since the curves $c^i$ all have the same projection to $M$ on $I$.
Then $C(0)=e$ and 
\[ \dot C(s)=\alpha_1\cdot_{TM}\dot c^1(s)+_{TM}\ldots+_{TM}\alpha_k\cdot_{TM} \dot c^k(s)
\]
by definition of $\cdot_{TM}$ and $+_{TM}$.
This yields 
\[ \dot C(s)=\alpha_1\cdot_{TM}\chi(c^1(s))+_{TM}\ldots+_{TM}\alpha_k\cdot_{TM} \chi(c^k(s))=\chi\left( \sum_{i=1}^k\alpha_i c^i(s)
\right)=\chi(C(s))
\]
for all $s\in I$,
since $\chi$ is linear. This shows that $C\colon I\to E$ is an integral curve of $\chi$ through $e$ at time $0$ and which is defined at time $t$. 
Hence $(t,e)\in \widetilde \Omega$. The above shows that 
\[ \{(t, e)\in \mathbb R\times E\mid (t, q(e))\in \Omega\}\subseteq \widetilde \Omega.
\]
Hence the two sets are equal.


\medskip
For each $t\in \R$ and each $x\in M_t$ the flow $\Phi_t$ restricts to a map
\[ E\an{x}\to E\an{\phi_t(x)}
\]
which is bijective since it has as inverse the restriction of $\Phi_{-t}$ to $E\an{\phi_t(x)}$.
To see that $\Phi_t$ is linear, note that 
\[ T\!\!+\circ(\chi,\chi)\colon E\times_ME\to TE
\]
equals 
\[ \chi\circ +\colon E\times_ME\to TE.
\]
Hence the vector field $(\chi, \chi)$ on $E\times_M E$, which is an embedded submanifold of $E\times E$, is $+$-related to the vector field $\chi$ on $E$. Since the flow $(\chi, \chi)$ at a time $t$ is $(\Phi_t,\Phi_t)$, this integrates to the equality
\[ +\circ (\Phi_t,\Phi_t)=\Phi_t\circ +,
\]
which exactly restricts to the additivity of $\Phi_t\an{E\an x}$. Similarly, for each $\alpha\in \R$, the compatibility of $\chi$ with the scalar multiplication by $\alpha$ integrates to the $\mathbb R$-homogeneity of $\Phi_t\an{E\an x}$, as already seen in the considerations above.
\end{proof}

\end{document}
